% 熱物理学 熱物理学-8 問題2（060-cv-t3） 解いた日 2026-11-16
%% kotae
\toi{1}{◎}{}
\[ \left(\pdv{S}{T}\right)_{V}=\frac{C_{V}}{T}=AT^{2} \]
\[ S(T,V)=\int AT^{2}\,\dd T=\frac{1}{3}AT^{3}+S_{0} \]
\[ \Delta S=\frac{1}{3}A\left(T_{2}^{3}-T_{1}^{3}\right) \]
\toi{2}{◎}{}
\[ \delta Q=\kakimachigai{\delta U}{\dd U}=C_{V}\,\dd T \]
\[ Q=\int_{T_{1}}^{T_{2}}AT^{3}\,\dd T=\frac{1}{4}A\left(T_{2}^{4}-T_{1}^{4}\right) \]
\toi{3}{○}{記述}
\begin{align*}
T_{\mathrm{m}}&=\frac{\frac{1}{4}A(T_{2}^{4}-T_{1}^{4})}{\frac{1}{3}A(T_{2}^{3}-T_{1}^{3})}\\
&=\frac{3(T_{2}^{4}-T_{1}^{4})}{4(T_{2}^{3}-T_{1}^{3})}\\
&\xrightarrow[T_{1}\to0]{}\frac{3}{4}T_{2}
\end{align*}
理想気体の $S$ の温度依存部分は $\ln(T_{2}/T_{1})$。$T_{1}\to0$ で $\ln(T_{2}/T_{1})\to\infty$ と発散する。\nuke{$\Delta S$ は有限}
\naze{問われている「$\Delta S$ が有限に留まる」こととその理由が無い。(1) から $T_{1}\to0$ で $\Delta S\to AT_{2}^{3}/3$ と有限であり，それは被積分関数 $C_{V}/T=AT^{2}$ が $T\to0$ で0になり積分できる（比熱が十分速く0に近づく）からである。理想気体の被積分関数は $nc_{V}/T$ なので発散する，と並べて述べる。}
%% kizuki
\item $\Delta S=\int C_{V}/T\,\dd T$ の下端 $T_{1}\to0$ での振る舞いは被積分関数で決まる。$C_{V}=AT^{3}$ なら $AT^{2}$ で積分は有限，比熱一定の理想気体なら $nc_{V}/T$ で $\ln T_{1}\to-\infty$ と発散する。2つを「被積分関数の違い」として並べて書く。
\item \mae{第3章 問題4 (5)・11-12} 前回は「比べよ」の比較が無く，今回は比較だけで本題（$\Delta S$ が有限）が無い。解き終えたら問題文の問い（「表せ」「述べよ」）を数え，それぞれに答えた式か文があるかを確かめる。
